% !TeX root = ../../Book2.tex
% 纲要标题：## 第17章　中心单代数
% 纲要位置：计划纲要.md，第 1594 行。

\chapter{中心单代数}
\label{b2:ch:17}
\BookChapterNumber{b2:ch:17}

\begin{chapterabstract}
中心单代数在适当扩张标量后成为矩阵代数，而它在原域上可能保留非交换除代数的结构。本章以正则双模和稠密性定理建立包络代数判据，证明任意标量扩张和张量积保持中心单性。中心化子定理给出简单子代数的维数关系，Skolem–Noether 定理比较它的不同嵌入。通过在除代数中寻找可分最大子域，建立有限 Galois 分裂域的存在；四元数的显式范数与分裂方程则为下一章的 Brauer 群提供可计算的例子。
\end{chapterabstract}

\begin{learninggoals}
\begin{itemize}
\item 区分单性、中心性与分裂性，证明中心单代数的包络代数判据及次数平方公式。
\item 证明标量扩张、下降与张量积结论，说明中心化子定理在子代数中心不可分时仍然成立。
\item 分清包含极大子域、最大交换子代数与分裂域，构造除代数的可分最大子域并证明有限 Galois 分裂域存在。
\item 用两个模结构证明 Skolem–Noether 定理，辨认其假设，并将四元数的分裂问题化为范数方程。
\end{itemize}
\end{learninggoals}

实四元数代数 \(\mathbb H\) 中，每个非零元素都有逆，因此不可能同构于 \(M_2(\mathbb R)\)：后者有非零而不可逆的矩阵。可是把标量扩到复数，\(\mathbb H\otimes_{\mathbb R}\mathbb C\) 却成为 \(M_2(\mathbb C)\)。同一个代数在原域上保留除法结构，在较大域上呈现为矩阵代数。由扩域后的形状倒推原代数，必须记录底域中的哪些信息仍可恢复。

\ProseParagraphBreak

没有非零真双边理想是单性，中心恰为指定底域是中心性，同构于底域上的全矩阵代数则是分裂性。三者要分开判断：非平凡有限域扩张作为基域代数是单的，中心却更大；实四元数在实数上中心单，仍不在实数上分裂。若其中嵌入一个简单子代数，它与自己的中心化子共同限制了整个代数的维数。进一步若要找到使代数分裂的扩域，就要比较子域、不同嵌入和标量扩张的作用；四元数提供了可以手算的检验场所。

\ProseParagraphBreak

本章需要第十一章的包络代数与双中心化子、第十三章的矩阵环分类和简单模重数，以及第十四章 Jacobson 稠密性定理的有限维结论。有关分裂域的论证使用第一卷定理15.2.3的代数闭包存在性，以及可分多项式和有限 Galois 扩张的基本性质。任意标量扩张保持中心单性不要求可分；有限可分分裂域的存在则需要另行构造可分最大子域。

\ProseParagraphBreak
中心单代数、中心化子与分裂域的系统论述可对照 Gille–Szamuely 的《Central Simple Algebras and Galois Cohomology》\cite[Chapter 2, pp. 19--69]{GilleSzamuely2017CentralSimpleAlgebras}。除代数的基础及实四元数可对照 Roman 的《Advanced Linear Algebra》\cite[Chapter 18, Division Algebras and The Quaternions, pp. 462--463; Theorem 18.10, p. 463]{Roman2008AdvancedLinearAlgebra}。关于有限维实除代数的进一步分类，该书\cite[Theorem 18.12, pp. 466--469]{Roman2008AdvancedLinearAlgebra}给出 Frobenius 定理。本章先建立适用于任意底域的中心单结构与嵌入理论，再在特征不为二时展开四元数的具体计算。

% !TeX root = ../../Book2.tex
% 纲要标题：### 17.1 单代数与中心
% 纲要位置：计划纲要.md，第 1596 行。

\section{单代数与中心}
\label{b2:sec:1701}

半单结构定理把有限维代数中的单块写成除环上的矩阵环，但底域未必就是这些块的中心。中心中的元素同全部乘法相容，所以每个单块也自然成为其中心上的代数。研究一个单块扩大标量后的结构时，就要先确定这个标量域：中心单代数要求它恰为全部中心，接下来将证明这一条件使任意扩域后的代数仍然单。张量积和中心化子的结构也将以此为基础。本章的代数结合且含单位元，子代数与嵌入保持单位元；中心单代数一律非零且有限维。

% 纲要要点（供目录核对）：
%
% * 有限维中心单代数；固定底域与实际中心，包络判据核对稠密性所用右除环及全线性算子目标
% * 分裂代数与分裂域；区别给定映射的同构检测与扩域后新同构的存在，核对有限维自同态的标量扩张、基比较下降及次数平方，正式证明纯不可分扩域后的非单性反例
% * 张量积与反代数；用两个包络作用证明张量积中心单，区别一般单代数的非单张量积
%

\subsection{有限维中心单代数}
\label{b2:subsec:1701:01}

\begin{definition}\label{b2:def:central-simple-algebra}
设 \(k\) 为域，\(A\) 为非零有限维含幺结合 \(k\) 代数。若 \(A\) 没有非零真双边理想，且其中心满足 \(Z(A)=k1_A\)，则称 \(A\) 为 \(k\) 上的\term[zhongxin-dan-daishu]{中心单代数}[central simple algebra]。下文用结构嵌入把 \(k1_A\) 识别为 \(k\)。
\end{definition}

单性限制的是双边理想，中心性还要求指定的底域就是全部中心。例如，有限域扩张 \(L/k\) 作为 \(k\) 代数是单的，但中心为 \(L\)，只有 \(L=k\) 时才中心单。对 \(n\ge1\)，矩阵代数 \(M_n(k)\) 则中心单：单性由\cref{b2:prop:simple-artinian}证明，中心为标量矩阵是\cref{b2:prop:matrix-centers-central-idempotents}在除环为 \(k\) 时的结论。

\ProseParagraphBreak
底域的选择确实会改变这个判断。例如 \(M_2(\mathbb C)\) 作为实代数是单的，中心却是 \(\mathbb C I_2\)，大于 \(\mathbb R I_2\)；它的实维数为 \(2\cdot2^2=8\)。改以 \(\mathbb C\) 为底域后，同一个乘法代数成为中心单代数，维数为 \(2^2=4\)。因此“中心单”要连同底域一起说明；后面得到的维数平方公式也只对这个指定的中心域成立。

\ProseParagraphBreak
一般非零含幺单代数 \(B\) 的中心是域，这不需要有限维假设。若 \(0\ne z\in Z(B)\)，则 \(Bz\) 是非零双边理想，故等于 \(B\)，从而 \(bz=1\) 对某个 \(b\) 成立。由 \(z\) 中心可知 \(zb=1\)，且从 \(zx=xz\) 可推出 \(z^{-1}x=xz^{-1}\)。所以 \(z\) 在中心内可逆。若 \(B\) 还是有限维 \(k\) 代数，\(Z(B)\) 作为 \(B\) 的 \(k\) 子空间也有限维，因而是 \(k\) 的有限扩张。

\ProseParagraphBreak
Artin–Wedderburn 定理给出 \(A\cong M_r(D)\) 时，除环 \(D\) 的中心未必就是原底域 \(k\)。中心单条件进一步固定这个参数，要求 \(Z(D)=k\)。

\begin{proposition}\label{b2:prop:csa-division-description}
设 \(k\) 为域，\(A\) 为中心单 \(k\) 代数，则存在整数 \(r\geq1\) 与有限维中心除 \(k\) 代数 \(D\)，使 \(A\cong M_r(D)\)。反过来，对任意这样的 \(r,D\)，矩阵代数 \(M_r(D)\) 中心单。整数 \(r\) 与 \(D\) 的 \(k\) 代数同构类型由 \(A\) 唯一确定。
\end{proposition}
\begin{proof}
有限维使左理想降链稳定，单性与\cref{b2:prop:simple-artinian}给出 \(A\cong M_r(D)\)。这个构造保持中心中的 \(k\) 标量，因此是 \(k\) 代数同构，\(D\) 有限维。由\cref{b2:prop:matrix-centers-central-idempotents}，\(Z\bigl(M_r(D)\bigr)=Z(D)I_r\)，而 \(A\) 的中心恰为 \(k\)，于是 \(Z(D)=k\)。反向的单性由\cref{b2:prop:simple-artinian}中矩阵环的结论得到，中心由同一中心公式给出。唯一性是\cref{b2:prop:semisimple-multiplicity}对正则模及其自同态除环的唯一性，构造同时保留 \(k\) 作用。
\end{proof}

\begin{theorem}[包络代数判据]\label{b2:thm:csa-enveloping-isomorphism}
设 \(k\) 为域，\(A\ne0\) 为有限维含幺结合 \(k\) 代数，则 \(A\) 中心单，当且仅当左右乘作用给出同构
\begin{equation}\label{b2:eq:csa-enveloping-isomorphism}
\Theta_A:A\otimes_kA^{\mathrm{op}}\longrightarrow\operatorname{End}_k(A),
\qquad a\otimes b^{\mathrm{op}}\longmapsto(x\mapsto axb).
\end{equation}
\end{theorem}
\begin{proof}
\cref{b2:thm:enveloping-bimodules}已证明 \(\Theta_A\) 是代数同态。若 \(A\) 中心单，要证明它满射，就需让左右乘的有限和实现任意 \(k\) 线性算子。把正则双模 \(A\) 看成左 \(A^e\) 模，其子模正是双边理想，故这个模简单。由\cref{b2:prop:regular-bimodule-endomorphism}，稠密性定理所用的右标量除环为
\[
D=\operatorname{End}_{A^e}(A)^{\mathrm{op}}=Z(A)^{\mathrm{op}}=k,
\]
所以目标 \(\operatorname{End}_D(A_D)\) 恰为 \(\operatorname{End}_k(A)\)。\(A\) 对这个右除环有限维，\cref{b2:cor:density-finite-dimensional}用于环 \(A^e\) 及简单模 \(A\)，便使 \(\Theta_A\) 满射。来源与目标的 \(k\) 维数同为 \((\dim_kA)^2\)，故它还是单射。

反过来，若 \(\Theta_A\) 同构，则每个双边理想 \(I\) 都是正则双模 \(A\) 的 \(A^e\) 子模，因而被 \(\Theta_A\) 的全部作用算子保持。满射性使它被每个 \(T\in\operatorname{End}_k(A)\) 保持。若 \(0\ne x\in I\)，对任意 \(y\in A\)，把 \(x\) 扩充为一组基并指定 \(T(x)=y\)，就得到 \(y=T(x)\in I\)。所以 \(I=A\)，证明单性。由正则双模自同态公式，\(Z(A)\) 同构于 \(\operatorname{End}_{A^e}(A)\)；后者现在是全部线性算子的中心化子。在\cref{b2:prop:tensor-factor-centralizers}中取 \(U=A\)、\(W=k\)，便知这个中心化子为 \(k\operatorname{id}_A\)，故 \(Z(A)=k\)。
\end{proof}

张量积中的元素是有限和，所以这个同构具体断言：对每个 \(T\in\operatorname{End}_k(A)\)，都能找到有限组 \(a_\nu,b_\nu\in A\)，使
\[
T(x)=\sum_{\nu=1}^{r}a_\nu x b_\nu\qquad(x\in A).
\]
纯张量 \(a\otimes b^{\mathrm{op}}\) 对应一次左右乘；一般线性算子由这样的操作的有限和实现。

\ProseParagraphBreak
在 \(A=M_n(k)\) 中，若 \(X=(x_{pq})\)，矩阵单位给出
\[
E_{ij}XE_{\ell m}=x_{j\ell}E_{im}.
\]
这项操作读取第 \((j,\ell)\) 个条目，再把它放到第 \((i,m)\) 个位置，其余条目清零。让两组位置遍历所有可能，就得到矩阵空间的全部线性自同态的矩阵单位；它们的线性组合正是任意线性操作。扩大底域后，这些读取、放置公式依然成立，这也显示了包络代数判据适合处理标量扩张的原因。

\subsection{分裂代数与分裂域}
\label{b2:subsec:1701:02}

对域扩张 \(K/k\)，记 \(A_K=K\otimes_kA\)。其乘法逐因子定义，单位为 \(1\otimes1\)。
\begin{definition}\label{b2:def:csa-splitting-field}
设 \(k\) 为域，\(A\) 为中心单 \(k\) 代数。若存在整数 \(n\geq1\) 使 \(A\cong M_n(k)\)，则称 \(A\) 为\term[fenlie-daishu]{分裂代数}[split algebra]，或称它在 \(k\) 上分裂。对域扩张 \(K/k\)，记 \(A_K=K\otimes_kA\)；若 \(A_K\cong M_n(K)\)，则称 \(K\) 为 \(A\) 的\term[fenlie-yu-daishu]{分裂域}[splitting field of an algebra]。
\end{definition}
\symidx{\(A_K=K\otimes_kA\)：代数的标量扩张}

扩张标量能检测一个给定映射是否为同构，但这与扩张后才找到一个同构不同。若 \(f:V\to W\) 是给定的 \(k\) 线性映射，且 \(K\otimes_k f\) 为同构，则 \(f\) 本身就是同构：向量空间的短正合列分裂，张量后仍正合，所以 \(f\) 的核与余核扩张后都为零；非零 \(k\) 向量空间的一组基扩张后仍非空，因而核与余核原本就为零。对于给定的 \(k\) 代数同态，这同样检测它是否为代数同构。

\ProseParagraphBreak
若只知道 \(A_K\cong B_K\)，这个同构却未必来自某个 \(k\) 代数同态 \(A\to B\)。例如实四元数代数 \(\mathbb H\) 与 \(M_2(\mathbb R)\) 扩张到 \(\mathbb C\) 后都同构于 \(M_2(\mathbb C)\)，其中四元数的分裂将在\cref{b2:prop:quaternion-basis-splitting}中具体给出；但 \(\mathbb H\) 是除代数，\(M_2(\mathbb R)\) 有非零不可逆矩阵，所以两者在 \(\mathbb R\) 上不同构。\cref{b2:cor:similarity-descent}中矩阵相似能够下降，是因为首一不变因子在扩张标量后保留，这个额外的分类信息不能直接推广到任意代数的同构问题。

\begin{theorem}[标量扩张]\label{b2:thm:csa-base-change}
设 \(k\) 为域，\(A\) 为中心单 \(k\) 代数，则对任意域扩张 \(K/k\)，标量扩张 \(A_K=K\otimes_kA\) 为中心单 \(K\) 代数，不要求扩张有限或可分。
\end{theorem}
\begin{proof}
要保留中心单性，可以把包络代数判据中的同构扩张标量。由于 \(A\) 有限维，取它的一组有限基后，矩阵单位给出自然同构
\[
K\otimes_k\operatorname{End}_k(A)\cong\operatorname{End}_K(A_K),
\]
将 \(c\otimes T\) 送到 \(d\otimes a\mapsto cd\otimes T(a)\)。它把两侧相应的有限矩阵单位基一一对应，所以是代数同构；这个与自同态代数交换的公式不能无条件用于无限维空间。另有
\[
A_K\otimes_K(A_K)^{\mathrm{op}}
\cong K\otimes_k(A\otimes_kA^{\mathrm{op}}).
\]
将\eqref{b2:eq:csa-enveloping-isomorphism}扩张标量并作这些识别，所得映射正是 \(A_K\) 的左右乘作用。它是同构，且 \(A_K\ne0\)，因为原 \(k\) 基扩张后仍为 \(K\) 基。因此包络代数判据给出所需中心单性。
\end{proof}

若只假设 \(A\) 单，标量扩张后却可能出现非零真双边理想。\cref{b2:prop:inseparable-base-change-counterexample}将给出一个有限纯不可分域扩张作为 \(k\) 代数的反例；它的中心大于 \(k\)，因而不满足上述定理的中心性条件。

\begin{proposition}[中心单性的下降]\label{b2:prop:csa-descent-from-extension}
设 \(k\) 为域，\(A\ne0\) 为有限维含幺结合 \(k\) 代数，\(K/k\) 为域扩张。若 \(A_K=K\otimes_kA\) 为中心单 \(K\) 代数，则 \(A\) 为中心单 \(k\) 代数。
\end{proposition}
\begin{proof}
若 \(I\) 是 \(A\) 的非零真双边理想，则 \(K\otimes_kI\) 是 \(A_K\) 的双边理想。取 \(I\) 的一组 \(k\) 基并扩充为 \(A\) 的基，便知扩张后它仍非零且仍为真子空间，这与 \(A_K\) 单矛盾。因此 \(A\) 单。

再把 \(1_A\) 扩充为 \(A\) 的 \(k\) 基。若 \(z\in Z(A)\)，则 \(1\otimes z\) 与 \(K\) 标量及全部 \(1\otimes a\) 交换，所以在 \(Z(A_K)=K(1\otimes1_A)\) 中。于是 \(1\otimes z=c\otimes1_A\) 对某个 \(c\in K\) 成立。用扩张后的基比较系数，\(z\) 的所有非单位基坐标都为零，其单位坐标原本在 \(k\) 中，故 \(z\in k1_A\)。这证明中心性。
\end{proof}

\begin{corollary}\label{b2:cor:csa-algebraic-closure-splits}
设 \(k\) 为域，\(\Omega\) 为其代数闭包，\(A\) 为中心单 \(k\) 代数，则 \(A_\Omega=\Omega\otimes_kA\cong M_n(\Omega)\)，其中整数 \(n\geq1\) 唯一，且 \(\dim_kA=n^2\)。
\end{corollary}
\begin{proof}
\cref{b2:thm:csa-base-change}使 \(A_\Omega\) 有限维且中心单，因而简单、半单。\cref{b2:prop:finite-dimensional-semisimple-algebra}在代数闭域上把它写成矩阵环的有限乘积；单性排除多个非零因子，所以只剩 \(M_n(\Omega)\)。扩张标量不改变有限基的大小，故 \(\dim_kA=n^2\)，这也唯一确定 \(n\)。
\end{proof}

称 \(n=\sqrt{\dim_kA}\) 为 \(A\) 的\term[zhongxin-dan-daishu-cishu]{次数}[degree of a central simple algebra]，记为 \(\deg A\)。次数为 \(n\)，整个代数的向量空间维数则为 \(n^2\)；因此中心单代数的维数必须是完全平方，例如三维或六维的代数不能中心单。任意分裂域给出的矩阵大小都为这个 \(n\)，因为维数在标量扩张下不变。\secref{b2:sec:1702}还将证明可以选取有限 Galois 分裂域；仅知道在代数闭包上分裂，尚未说明有限可分扩张的存在。
\symidx{\(\deg A=\sqrt{\dim_kA}\)：中心单代数的次数}

\subsection{张量积与反代数}
\label{b2:subsec:1701:03}

反代数 \(A^{\mathrm{op}}\) 与 \(A\) 有相同的中心，双边理想也相同，所以中心单性在取反代数后保持。它的次数等于 \(\deg A\)，因为维数不变。与自身的反代数作张量则有更强的结论：由\cref{b2:thm:csa-enveloping-isomorphism}，
\[
A\otimes_kA^{\mathrm{op}}\cong M_{\dim_kA}(k)
=M_{(\deg A)^2}(k).
\]
无论 \(A\) 本身是否分裂，这个张量积总是分裂的。右侧矩阵的阶数是 \(\dim_kA=(\deg A)^2\)，不能误记为 \(\deg A\)。在\secref{b2:sec:1801}的 Brauer 群中，分裂矩阵代数代表单位元，这个同构将使 \(A^{\mathrm{op}}\) 的类成为 \(A\) 的类的逆元。

\begin{theorem}[中心单代数的张量积]\label{b2:thm:csa-tensor-product}
设 \(k\) 为域，\(A,B\) 为中心单 \(k\) 代数，则 \(A\otimes_kB\) 也是中心单 \(k\) 代数，且
\(\deg(A\otimes_kB)=\deg A\deg B\)。
\end{theorem}
\begin{proof}
令 \(C=A\otimes_kB\)。\(C\) 的左右乘可以分别在 \(A,B\) 两个因子上实现，因此要把它的包络作用与两个因子的包络作用相比较。交换不同张量因子的次序并取结合识别，得到
\[
C\otimes_kC^{\mathrm{op}}
\cong(A\otimes_kA^{\mathrm{op}})\otimes_k(B\otimes_kB^{\mathrm{op}}).
\]
分别应用两个包络代数同构，再用\cref{b2:prop:tensor-factor-centralizers}中已经通过矩阵单位证明的同构，得到
\[
C\otimes_kC^{\mathrm{op}}
\cong\operatorname{End}_k(A)\otimes_k\operatorname{End}_k(B)
\cong\operatorname{End}_k(C).
\]
在纯张量上，这个复合把 \((a\otimes b)\otimes(a'\otimes b')^{\mathrm{op}}\) 送到
\(x\otimes y\mapsto axa'\otimes byb'\)，正是 \(C\) 的左右乘作用。因此包络代数判据保证 \(C\) 中心单。维数乘法给出次数乘法。
\end{proof}

特别地，\(M_r(A)\cong M_r(k)\otimes_kA\) 是中心单代数，次数为 \(r\deg A\)。同构把 \(E_{ij}\otimes a\) 送到第 \((i,j)\) 个条目为 \(a\)、其余为零的矩阵；基与乘法同时对应。另一方面，两个一般单 \(k\) 代数的张量积不必单，中心性在定理中不能删除。例如
\[
\mathbb C\otimes_{\mathbb R}\mathbb C\longrightarrow\mathbb C\times\mathbb C,
\qquad z\otimes w\longmapsto(zw,z\bar w)
\]
是实代数同构。确实，按第一个 \(\mathbb C\) 因子看，来源以 \(1\otimes1,1\otimes i\) 为基，两个像为 \((1,1),(i,-i)\)，线性无关且生成目标；乘法直接保持。恒等嵌入和复共轭使第二因子分别产生两个坐标，而目标中的 \(\mathbb C\times0\) 是非零真双边理想，故张量积不是单环。这里每个因子作为实代数都单，但其中心为 \(\mathbb C\)，大于底域 \(\mathbb R\)。

\ProseParagraphBreak
把底域换成交换环时，有限维性应改为忠实有限投射性；左右乘作用的同构则仍有意义。\appref{b2:app:C}的\cref{b2:def:azumaya-algebra}据此定义 Azumaya 代数\footnote{東屋五郎（假名：あずまや ごろう，罗马音：Gorō Azumaya，1920-2010），日本数学家，研究环论及代数的结构；Azumaya 代数以他命名。}，并由\cref{b2:thm:azumaya-fiber-criterion}证明其剩余域纤维判据。这项推广不将每个纤维预设为已经分裂的矩阵代数。

\begin{proposition}[不可分标量扩张后的非单性]\label{b2:prop:inseparable-base-change-counterexample}
设 \(p\) 为素数、\(t\) 为 \(\mathbb F_p\) 上的超越元，令 \(k=\mathbb F_p(t)\)。在 \(k\) 的代数闭包中取 \(\alpha\) 满足 \(\alpha^p=t\)，并令 \(L=k(\alpha)\)。则 \(X^p-t\) 在 \(k[X]\) 中不可约，\([L:k]=p\)，且以第一个张量因子为 \(L\) 标量时有含幺 \(L\) 代数同构
\[
L\otimes_kL\cong L[\varepsilon]/(\varepsilon^p).
\]
其中 \(1\otimes\alpha-\alpha\otimes1\) 对应 \(\varepsilon\) 的剩余类。它是非零幂零元，生成一个非零真理想，故 \(L\otimes_kL\) 不是单代数。原代数 \(L\) 虽为单 \(k\) 代数，但其中心为 \(L\ne k\)。
\end{proposition}
\begin{proof}
先证明 \(t\) 在 \(k\) 中不是 \(p\) 次方。若 \((f/g)^p=t\)，其中 \(f,g\in\mathbb F_p[t]\)、\(g\ne0\)，约去公因子后有 \(f^p=tg^p\)。左边不可约因子 \(t\) 的指数为 \(p\) 的倍数，右边的指数则模 \(p\) 余一，矛盾。

令 \(\mu_{\alpha,k}\) 为 \(\alpha\) 在 \(k\) 上的首一最小多项式。它整除 \(X^p-t\)，而在代数闭包中 \(X^p-t=(X-\alpha)^p\)，故 \(\mu_{\alpha,k}=(X-\alpha)^d\)，其中 \(1\le d\le p\)。若 \(d<p\)，其 \(X^{d-1}\) 系数为 \(-d\alpha\in k\)。此时 \(d\) 在特征 \(p\) 的域中非零，便有 \(\alpha\in k\)，与 \(t\) 不是 \(p\) 次方矛盾。因此 \(d=p\)，\(\mu_{\alpha,k}=X^p-t\)，并且 \(1,\alpha,\ldots,\alpha^{p-1}\) 是 \(L/k\) 的基。

把 \(L[X]\) 中的 \(X\) 送到 \(1\otimes\alpha\)，把系数 \(c\in L\) 送到 \(c\otimes1\)，得到含幺 \(L\) 代数同态。关系 \(X^p-t\) 在像中为零，故它诱导
\[
L[X]/(X^p-t)\longrightarrow L\otimes_kL.
\]
两侧分别以 \(X^j\) 的剩余类和 \(1\otimes\alpha^j\)（\(0\le j<p\)）为 \(L\) 基，且这些基向量一一对应，所以此映射为同构。在 \(L[X]\) 中令 \(\varepsilon=X-\alpha\)，便有 \(X^p-t=\varepsilon^p\)，得到所述截断多项式环。它以 \(1,\varepsilon,\ldots,\varepsilon^{p-1}\) 的剩余类为 \(L\) 基，故 \(\varepsilon\ne0\) 而 \(\varepsilon^p=0\)；理想 \((\varepsilon)\) 非零，且商为非零域 \(L\)，所以是真理想。最后，\(L\) 是域，故作为 \(k\) 代数单；但 \([L:k]=p>1\)，且 \(Z(L)=L\)，因而它不中心单，也就不满足\cref{b2:thm:csa-base-change}的假设。
\end{proof}

% !TeX root = ../../Book2.tex
% 纲要标题：### 17.2 中心化子、子域与分裂域
% 纲要位置：计划纲要.md，第 1313 行。

\section{中心化子、子域与分裂域}
\label{b2:sec:1702}

一个单子代数 \(B\subseteq A\) 的中心化子记录与它交换的另一部分乘法。维数公式将精确说明这两部分各占多大；但只有 \(B\) 的中心也是 \(k\) 时，二者才直接张量成整个 \(A\)。本节通过简单模及其重数证明这些结论，子代数中心的扩张 \(Z(B)/k\) 可以不可分。

% 纲要要点（供目录核对）：
%
% * 中心化子的维数关系
% * 双中心化子定理
% * 子域与最大交换子代数
% * 可分最大子域与有限 Galois 分裂域
%

\subsection{中心化子的维数关系}
\label{b2:subsec:1702:01}

记
\[
C_A(B)=\{\,a\in A\mid ab=ba\text{ 对所有 }b\in B\,\}.
\]
它是含单位元的 \(k\) 子代数。本节的中心化子在 \(A\) 内取，与第十一章在整个线性自同态代数内取中心化子应分别理解。
\symidx{\(C_A(B)\)：子代数 \(B\) 在 \(A\) 中的中心化子}

\begin{theorem}[中心化子定理]\label{b2:thm:csa-centralizer}
设 \(k\) 为域，\(A\) 为中心单 \(k\) 代数，\(B\subseteq A\) 为非零简单含幺 \(k\) 子代数，\(L=Z(B)\)。令 \(C=C_A(B)=\{\,a\in A\mid ab=ba\text{ 对所有 }b\in B\,\}\)，则
\[
C\text{ 为简单代数},\qquad Z(C)=L,
\qquad \dim_kA=\dim_kB\,\dim_kC,
\qquad C_A(C)=B.
\]
其中两个中心均按它们在 \(A\) 中的实际嵌入识别。结论不要求 \(L/k\) 可分。
\end{theorem}
\begin{proof}
先构造有限维代数 \(E=B\otimes_kA^{\mathrm{op}}\)。\(B\) 以其中心 \(L\) 为底域是中心单代数，而
\[
E\cong B\otimes_L(L\otimes_kA^{\mathrm{op}}).
\]
具体同构将 \(b\otimes(l\otimes a^{\mathrm{op}})\) 送到 \(bl\otimes a^{\mathrm{op}}\)，逆映射在 \(b\otimes a^{\mathrm{op}}\) 上取 \(b\otimes(1\otimes a^{\mathrm{op}})\)。\cref{b2:thm:csa-base-change}说明第二个因子在 \(L\) 上中心单，再由\cref{b2:thm:csa-tensor-product}，\(E\) 为中心单 \(L\) 代数，特别是单 Artin 环。

令 \(E\) 在 \(A\) 上作用为 \((b\otimes a^{\mathrm{op}})x=bxa\)。一个 \(E\) 线性自同态首先与全部右乘交换，所以由\cref{b2:thm:regular-double-centralizer}，它形如左乘 \(c\in A\)。再要求它与左 \(B\) 作用交换，恰好得到 \(c\in C\)。左乘保留乘法次序，因此
\[
C\cong\operatorname{End}_E(A).
\]
写 \(E\cong M_r(D)\)，其中 \(D\) 是除环。\cref{b2:thm:artin-wedderburn}及\cref{b2:prop:semisimple-multiplicity}的矩阵块分类说明，\(E\) 半单，且简单左模只有列模 \(D^r\) 一种类型。\(A\) 是有限维 \(k\) 空间，所以作为 \(E\) 模为有限个副本之和，设 \(A\cong(D^r)^m\)、\(m\ge1\)。因此
\[
\operatorname{End}_E(A)\cong M_m(D^{\mathrm{op}}),
\]
从而 \(C\) 简单。

记 \(d=\dim_kD\)。分别比较上述代数与模的 \(k\) 维数，有
\[
\dim_kB\,\dim_kA=r^2d,\qquad
\dim_kA=mrd,\qquad\dim_kC=m^2d.
\]
这三式给出 \(\dim_kB\,\dim_kC=\dim_kA\)。至此已独立证明：对任意满足定理假设的简单子代数，其中心化子简单，且有上述维数公式；双中心化子和中心的识别尚未用到。现在 \(C\) 也是 \(A\) 的简单含幺子代数，将这个已证的中间结论用于 \(C\subseteq A\)，得到
\(\dim_k C_A(C)=\dim_kA/\dim_kC=\dim_kB\)。而 \(B\subseteq C_A(C)\) 由交换性定义直接成立，有限维与等维使二者相等。最后
\[
Z(C)=C\cap C_A(C)=C\cap B=Z(B)=L,
\]
得到中心的实际识别。
\end{proof}

证明中 \(L\otimes_kA^{\mathrm{op}}\) 的中心单性来自对任意域扩张均成立的标量扩张定理。即使 \(L/k\) 是纯不可分扩张，这一步也仍然有效；真正需要简单性的，是 \(B\) 以其中心为底域的结构。

\subsection{双中心化子定理}
\label{b2:subsec:1702:02}

\begin{corollary}\label{b2:cor:csa-centralizer-factorization}
设 \(k\) 为域，\(A\) 为中心单 \(k\) 代数，\(B\subseteq A\) 为非零简单含幺 \(k\) 子代数，记 \(C_A(B)\) 为 \(B\) 在 \(A\) 中的中心化子。若 \(Z(B)=k\)，则乘法诱导中心单 \(k\) 代数同构
\[
B\otimes_kC_A(B)\longrightarrow A,\qquad b\otimes c\longmapsto bc.
\]
一般地，令 \(L=Z(B)\)，则有 \(L\) 代数同构
\[
B\otimes_LC_A(B)\cong C_A(L).
\]
\end{corollary}
\begin{proof}
两个因子的元素互相交换，所以乘法公式保持张量关系与代数乘法。若 \(L=k\)，来源由\cref{b2:thm:csa-tensor-product}为单代数；映射保持单位元，故核不能为全代数，只能为零。中心化子维数公式说明来源与目标同维，于是满射。

一般情形中，\(B,C_A(B)\) 都是中心单 \(L\) 代数，故它们在 \(L\) 上的张量积单，乘法仍单射，且其像包含于 \(C_A(L)\)。这里仍以 \(k\) 为底域，将简单 \(k\) 子代数 \(L\subseteq A\) 代入中心化子定理，得到 \(\dim_k C_A(L)=\dim_kA/[L:k]\)。张量积来源的 \(k\) 维数为
\(\dim_kB\,\dim_kC_A(B)/[L:k]\)，也等于这个数，因此像就是全部 \(C_A(L)\)。
\end{proof}

例如，把有限扩张 \(L/k\) 通过正则作用嵌入 \(A=\operatorname{End}_k(L)\)。与左乘 \(L\) 交换的映射恰为右乘 \(L\)，因为 \(T(x)=T(x\cdot1)=x\cdot T(1)\)。\(L\) 交换，所以 \(C_A(L)=L\)，维数公式给出 \([L:k]^2=[L:k]\,[L:k]\)。当 \(L\ne k\) 时，\(L\otimes_L L=L\) 远小于 \(A\)；只有在正确的中心上取张量，且把目标写成 \(C_A(L)\)，才能保持一般公式成立。

\ProseParagraphBreak
双中心化子等式还说明，简单子代数可由它在 \(A\) 中的全部交换关系恢复。若删除简单性，这一性质可能失败：第十一章的上三角代数 \(T_2(k)\subseteq M_2(k)\) 的中心化子为 \(kI\)，再取一次得到整个 \(M_2(k)\)，并非原上三角代数。

这些关系可以在 \(A\) 内按包含画成两支：
\[
\begin{tikzcd}[column sep=2.2em,row sep=1.5em]
& & B \arrow[dr,hook] & & \\
k \arrow[r,hook] & L=Z(B)=Z(C) \arrow[ur,hook] \arrow[dr,hook]
& & C_A(L) \arrow[r,hook] & A\\
& & C=C_A(B) \arrow[ur,hook] & &
\end{tikzcd}
\]
两支相乘得到 \(B\otimes_LC\cong C_A(L)\)，而 \(C_A(C)=B\)、\(\dim_kA=\dim_kB\,\dim_kC\)。当 \(L=k\) 时，图中的 \(C_A(L)\) 才是整个 \(A\)。

\subsection{子域与最大交换子代数}
\label{b2:subsec:1702:03}

设 \(A\) 为中心单 \(k\) 代数。若含 \(k\) 的子域 \(L\subseteq A\) 不包含于更大的 \(k\) 子域，称 \(L\) 为\term[baohan-jida-ziyu]{包含极大子域}[inclusion-maximal subfield]。若一个含幺交换子代数不包含于更大的含幺交换子代数，则称它为\term[zuida-jiaohuan-zidaishu]{最大交换子代数}[maximal commutative subalgebra]。前者只在域中比较，后者也允许零因子或幂零元。文献中的 maximal subfield 还可能专指满足 \([L:k]=\deg A\) 的子域；本书在一般中心单代数中分别写明这两个条件，只在除代数语境把包含极大子域简称\term[zuida-ziyu]{最大子域}[maximal subfield]。

\cref{b2:tab:csa-subfield-conditions}按对象所在位置列出容易混淆的条件。最后一行的 \(K\) 只须是 \(k\) 的扩域，可以完全不嵌入 \(A\)。
\begin{table}[htbp]
\centering
\caption{中心单代数中的子域、交换子代数与分裂域}\label{b2:tab:csa-subfield-conditions}
\begin{tabularx}{\linewidth}{@{}p{0.31\linewidth}X@{}}
\toprule
对象 & 条件 \\
\midrule
包含极大子域 \(L\subseteq A\) & 不包含于 \(A\) 中更大的含 \(k\) 子域。\\
最大交换子代数 \(E\subseteq A\) & 不包含于 \(A\) 中更大的含幺交换 \(k\) 子代数。\\
自中心化子域 \(L\subseteq A\) & \(C_A(L)=L\)。\\
次数达到 \(\deg A=n\) 的子域 \(L\subseteq A\) & \([L:k]=n\)。\\
分裂域 \(K/k\) & \(A_K\cong M_n(K)\)；不要求 \(K\subseteq A\)。\\
\bottomrule
\end{tabularx}
\end{table}

\begin{proposition}\label{b2:prop:csa-self-centralizing-field}
设 \(k\) 为域，\(A\) 为中心单 \(k\) 代数，\(n=\deg A\)，\(L\subseteq A\) 为含 \(k\) 的子域，则以下条件等价：
\[
[L:k]=n;\qquad C_A(L)=L;\qquad
L\text{ 是 }A\text{ 的最大交换子代数}.
\]
若 \(A\) 本身为除代数，则每个含 \(k\) 的最大子域都满足这些条件。
\end{proposition}
\begin{proof}
将\cref{b2:thm:csa-centralizer}用于简单子代数 \(L\)，得到 \(n^2=[L:k]\dim_k C_A(L)\)。由于 \(L\subseteq C_A(L)\)，当 \([L:k]=n\) 时二者同维，所以相等；反向相等也立即给出次数条件。

若 \(C_A(L)=L\)，任何包含 \(L\) 的交换子代数都包含于该中心化子，因此只能是 \(L\)。反过来，若 \(c\in C_A(L)\)，则 \(L[c]\) 是包含 \(L\) 的交换子代数，最大交换性使 \(c\in L\)。最后，若环境为除代数 \(D\)，则 \(L[c]\) 是有限维交换整环，因而为域：非零元的乘法单射，有限维使其满射，从而有逆。最大子域的性质就迫使 \(L[c]=L\)，也得到中心化子等于 \(L\)。
\end{proof}

在矩阵代数中，包含极大不保证次数达到 \(\deg A\)。例如 \(M_2(\mathbb C)\) 中每个有限维的 \(\mathbb C\) 子域都等于 \(\mathbb C\)，所以标量域是包含极大子域，却不是最大交换子代数，也只有次数一。对角代数 \(\mathbb C\times\mathbb C\) 是最大交换子代数，但不是域；\(\mathbb C[N]\)、\(N=\begin{pmatrix}0&1\\0&0\end{pmatrix}\)，也是最大交换子代数，因为与 \(N\) 交换的矩阵恰为 \(aI+bN\)，而它含有非零幂零元。

\begin{proposition}\label{b2:prop:maximal-field-splits-csa}
设 \(k\) 为域，\(A\) 为中心单 \(k\) 代数，\(L\subseteq A\) 为含 \(k\) 的子域。若 \([L:k]=\deg A=n\)，则 \(L\) 分裂 \(A\)。特别地，每个有限维中心除 \(k\) 代数的最大子域都是它的分裂域。
\end{proposition}
\begin{proof}
把 \(A\) 看成右 \(L\) 向量空间，其维数为 \(\dim_kA/[L:k]=n\)。作用
\[
L\otimes_kA\longrightarrow\operatorname{End}_L(A),
\qquad l\otimes a\longmapsto(x\mapsto axl)
\]
是 \(L\) 代数同态；右 \(L\) 线性用到 \(L\) 内元素交换，乘法相容性由结合律与相同交换性给出。目标中的 \(A\) 始终是原来的右 \(L\) 空间。来源由\cref{b2:thm:csa-base-change}为单代数，所以保持单位元的映射单射。两侧的 \(L\) 维数均为 \(n^2\)，故为同构，目标为 \(M_n(L)\)。除代数的最大子域由\cref{b2:prop:csa-self-centralizing-field}满足 \([L:k]=\deg A\)，所以也适用。
\end{proof}

\subsection{可分最大子域与有限 Galois 分裂域}
\label{b2:subsec:1702:04}

有限维除代数确有最大子域：子域维数是有界正整数，取具有最大维数的一个即可。不过，这样得到的子域未必可分。为得到有限 Galois 分裂域，下面另作一次可分元素的选择。

\begin{theorem}[有限可分分裂域的存在]\label{b2:thm:csa-separable-splitting}
设 \(k\) 为域。每个中心单 \(k\) 代数都有有限可分分裂域，因而也有有限 Galois 分裂域。更具体地，每个有限维中心除 \(k\) 代数 \(D\) 都含有一个可分最大子域 \(L\)，满足 \([L:k]=\deg D\)，且 \(L\) 分裂 \(D\)。
\end{theorem}
\begin{proof}
由\cref{b2:prop:csa-division-description}，只须对中心除代数 \(D\) 找到所述 \(L\)：一旦 \(L\) 分裂 \(D\)，它也分裂 \(M_r(D)\)。若 \(k\) 有限，则 \(D\) 的元素总数也有限，由\cref{b2:thm:wedderburn-little}，\(D\) 交换；中心为 \(k\) 便迫使 \(D=k\)，取 \(L=k\) 即可。以下设 \(k\) 无限，令 \(n=\deg D\)，\(\Omega\) 为代数闭包。

取同构 \(D_\Omega\cong M_n(\Omega)\) 及 \(D\) 的 \(k\) 基 \(e_1,\ldots,e_{n^2}\)，其像为矩阵 \(E_1,\ldots,E_{n^2}\)。这些矩阵是 \(M_n(\Omega)\) 的 \(\Omega\) 基。要寻找 \(a=\sum_i c_ie_i\)、\(c_i\in k\)，使对应矩阵具有 \(n\) 个互异特征值。

先建立一个明确的多项式检验。对任意 \(H\in M_n(\Omega)\)，令 \(f_H(t)=\det(tI-H)\)，并令 \(U_j\) 表示次数小于 \(j\) 的 \(\Omega\) 多项式空间，\(U_0=0\)。考虑等维空间之间的映射
\[
U_{n-1}\oplus U_n\longrightarrow U_{2n-1},
\qquad(u,v)\longmapsto u f_H+v f_H'.
\]
在单项式基下取其行列式 \(\Delta(H)\)。这是首一多项式 \(f_H\) 与 \(f_H'\) 的结式 \(\operatorname{Res}(f_H,f_H')\) 的一种归一化，也可视为判别式；更换这些空间的固定基只会使行列式乘以一个非零常数。矩阵条目来自 \(f_H\) 及其导数的系数，所以 \(\Delta(H)\) 是 \(H\) 条目的多项式。

该行列式非零当且仅当 \(f_H,f_H'\) 互素。如果互素且 \(u f_H+v f_H'=0\)，则 \(f_H\mid v\)，但 \(\deg v<n\)，故 \(v=0\)，继而 \(u=0\)，映射单射而可逆。若有正次数公因子 \(h\)，则 \((u,v)=(f_H'/h,-f_H/h)\) 是符合次数限制的非零核向量，所以行列式为零。互素又等价于 \(f_H\) 没有重根：根 \(\lambda\) 同时使 \(f_H,f_H'\) 为零，恰好是它在因式分解中出现至少两次。

于是
\[
P(X_1,\ldots,X_{n^2})=\Delta\left(\sum_iX_iE_i\right)
\in\Omega[X_1,\ldots,X_{n^2}]
\]
是非零多项式。因为 \(E_i\) 张成全部矩阵，可在 \(\Omega\) 中选择系数使其和为具有 \(n\) 个互异对角元的对角矩阵，此处 \(\Delta\ne0\)。

即使 \(P\) 的系数在 \(\Omega\) 中，它也不能在无限子域 \(k\) 的全部点上消失。证明可对变量数归纳：一元非零多项式只有有限多个根；多元时先将它按最后一个变量展开，选一个非零系数多项式，用归纳假设为前面的变量取 \(k\) 中的值，使该系数非零，再用一元结论选最后一个值。因此可取 \(c_i\in k\) 使 \(P(c_1,\ldots,c_{n^2})\ne0\)。相应的 \(a=\sum_i c_ie_i\in D\) 在分裂矩阵模型中有 \(n\) 个互异特征值。任意零化矩阵的多项式必须在这 \(n\) 个特征值处为零，故次数至少为 \(n\)；其 \(n\) 次特征多项式又由\cref{b2:thm:exterior-cayley-hamilton}零化该矩阵。因此矩阵最小多项式就是这个没有重根的特征多项式。

\(a\) 在 \(k\) 上的最小多项式也正是上述矩阵的最小多项式。事实上，对左乘算子 \(L_a:D\to D\)，有 \(p(L_a)=0\) 当且仅当 \(p(a)=0\)：必要性只需将算子作用于 \(1\)。\cref{b2:prop:scalar-extension-invariant-factors}说明 \(L_a\) 扩张到 \(\Omega\) 后最小多项式原样不变；在 \(D_\Omega\cong M_n(\Omega)\) 上，这个扩张算子就是矩阵 \(H\) 的左乘，其最小多项式又等于 \(H\) 的最小多项式。因此 \(a\) 的首一最小多项式次数为 \(n\)，且在 \(\Omega\) 中无重根，所以这个最小多项式可分。这里没有预先假定 \(H\) 的特征多项式系数属于 \(k\)。

\(D\) 为除环，使交换子代数 \(k[a]\) 是有限维整环，因而为域。于是 \(L=k[a]\) 是次数 \(n\) 的可分子域，由\cref{b2:prop:csa-self-centralizing-field}及\cref{b2:prop:maximal-field-splits-csa}，它是最大子域并分裂 \(D\)。最后把 \(a\) 送到其最小多项式在 \(\Omega\) 中的一个根，将 \(L\) 嵌入 \(\Omega\)，并取该多项式在 \(\Omega\) 中的分裂域 \(F\)。第一卷命题15.1.2保证其有限性；它由这个可分多项式的全部根生成，按该卷定义19.1.1后的分裂域刻画，是有限 Galois 扩张，并且含有所选的 \(L\) 的像。继续扩张分裂矩阵环仍为矩阵环，所以 \(F\) 也分裂原来的 \(A\)。
\end{proof}

% !TeX root = ../../Book2.tex
% 纲要标题：### 17.3 Skolem–Noether 定理
% 纲要位置：计划纲要.md，第 1609 行。

\section{Skolem–Noether 定理}
\label{b2:sec:1703}

中心化子定理比较一个子代数与同它交换的部分。现在比较同一个简单代数嵌入 \(A\) 的两种方式：结论是，它们只相差 \(A\) 中一个可逆元的内共轭。证明不从方程组中猜这个可逆元，而是把两种嵌入看成两个模结构，再由简单 Artin 环的模分类构造所需同构。

% 纲要要点（供目录核对）：
%
% * 单子代数的嵌入及内自同构；两种嵌入只改变左作用，单 Artin 环的一种简单模类型使维数决定重数
% * 中心单代数的自同构；从右正则模同构在单位元的像构造共轭元，逐个自同构可提升不保证存在群截面
% * 适用假设与非单代数的反例；非单源的不同分量重数及非中心环境的外自同构分别说明失败原因
%

\subsection{单子代数的嵌入与内自同构}
\label{b2:subsec:1703:01}

\begin{theorem}[Skolem–Noether]\label{b2:thm:skolem-noether}
设 \(k\) 为域，\(A\) 为中心单 \(k\) 代数，\(B\) 为非零有限维简单含幺 \(k\) 代数。对任意两个保幺 \(k\) 代数同态 \(\varphi,\psi:B\to A\)，存在 \(u\in A^\times\)，使
\[
\psi(b)=u\cdot\varphi(b)u^{-1}\qquad(b\in B).
\]
这些同态因 \(B\) 简单而自动单射；\(Z(B)/k\) 不必可分。这个版本也见\cite[Main Theorem 7.7.1]{Voight2021QuaternionAlgebras}。
\end{theorem}
\begin{proof}
把 \(B\) 的左作用与 \(A\) 的右正则作用放在同一个代数 \(E=B\otimes_kA^{\mathrm{op}}\) 中，便能只改变嵌入给出的左作用而保留右作用。令 \(L=Z(B)\)。与\cref{b2:thm:csa-centralizer}的证明相同，
\[
E\cong B\otimes_L(L\otimes_kA^{\mathrm{op}})
\]
是两个中心单 \(L\) 代数的张量积，因此为单 Artin 环。这一步只用任意标量扩张的中心单性，不要求 \(L/k\) 可分。

在同一个 \(k\) 向量空间 \(A\) 上定义两种左 \(E\) 模结构：
\[
(b\otimes a^{\mathrm{op}})\cdot_\varphi x=\varphi(b)xa,
\qquad
(b\otimes a^{\mathrm{op}})\cdot_\psi x=\psi(b)xa.
\]
左右乘法交换，所以模公理成立。由\cref{b2:prop:simple-artinian}写 \(E\cong M_r(D)\)，再由\cref{b2:thm:artin-wedderburn}及\cref{b2:prop:semisimple-multiplicity}的矩阵块分类，简单左模只有 \(S=D^r\) 一种类型。两种模都是有限个 \(S\) 的直和，且底层 \(k\) 空间都是 \(A\)，因此
\[
A_\varphi\cong S^{\oplus m}\cong A_\psi,\qquad
m=\frac{\dim_kA}{\dim_kS}.
\]
这里可以只凭维数确定重数，正因为 \(E\) 是单 Artin 环，只有一种简单左模。一般半单环可能有多种简单类型，它们的重数不能只由总维数恢复。于是存在 \(E\) 模同构 \(F:A_\varphi\to A_\psi\)。

\(E\) 线性首先保证 \(F(xa)=F(x)a\)，即 \(F\) 是右正则模之间的同构。右正则模由 \(1\) 生成，因此令 \(u=F(1)\)，便有 \(F(x)=F(1\cdot x)=ux\)。逆同构也由一的像决定，故为左乘某个 \(v\in A\)；两次复合为恒等在 \(1\) 处给出 \(vu=uv=1\)，所以 \(u\) 可逆。再由 \(B\) 的作用相容性，
\[
u\cdot\varphi(b)=F\bigl(\varphi(b)\bigr)=\psi(b)F(1)=\psi(b)u.
\]
右乘 \(u^{-1}\) 得到所需公式。最后，同态的核是 \(B\) 的双边理想，且因保持单位元不能为 \(B\)，所以核为零。
\end{proof}

这称为\term[Skolem-Noether-dingli]{Skolem–Noether 定理}[Skolem--Noether theorem]。若 \(B_1,B_2\subseteq A\) 为简单 \(k\) 子代数，且给定同构 \(\alpha:B_1\to B_2\)，把 \(B_1\) 的包含与 \(\alpha\) 后接 \(B_2\) 的包含作为两种嵌入，就得到 \(\alpha\) 可延伸为 \(A\) 的一个内自同构。所选 \(u\) 一般不唯一；若 \(u,v\) 实现同一嵌入之间的共轭，则 \(v^{-1}u\) 恰好与 \(\varphi(B)\) 的全部元素交换。

\subsection{中心单代数的自同构}
\label{b2:subsec:1703:02}

\begin{corollary}\label{b2:cor:csa-automorphisms-inner}
设 \(k\) 为域，\(A\) 为中心单 \(k\) 代数，则 \(A\) 的每个 \(k\) 代数自同构都是内自同构，并有群同构
\[
\operatorname{Aut}_{k\text{-alg}}(A)\cong A^\times/k^\times.
\]
特别地，对每个整数 \(n\ge1\)，\(\operatorname{Aut}_{k\text{-alg}}\bigl(M_n(k)\bigr)\cong\operatorname{PGL}_n(k)\)，其中
\(\operatorname{PGL}_n(k)=\operatorname{GL}_n(k)/k^\times I_n\)。
\end{corollary}
\begin{proof}
在\cref{b2:thm:skolem-noether}中取 \(B=A\)，两个嵌入分别为恒等映射和给定自同构，得到满群同态
\(A^\times\to\operatorname{Aut}_{k\text{-alg}}(A)\)、\(u\mapsto(a\mapsto uau^{-1})\)。它的核由与全部 \(A\) 交换的可逆元组成，即 \(Z(A)^\times=k^\times\)。因此把商类 \(uk^\times\) 送到该共轭自同构，给出所述商群同构。矩阵代数情形直接代入。
\end{proof}
\symidx{\(\operatorname{PGL}_n(k)\)：一般线性群模去非零标量矩阵所得的商群}

例如，若 \(\operatorname{char}k\ne2\)，\(d\in k^\times\) 非平方，矩阵
\[
T=\begin{pmatrix}0&d\\1&0\end{pmatrix}
\]
给出二次域 \(k(\sqrt d)\) 的嵌入。其非平凡自同构 \(\sqrt d\mapsto-\sqrt d\) 由 \(P=\operatorname{diag}(1,-1)\) 实现，因为 \(PTP^{-1}=-T\)。这个式子直接显示了域自同构如何在更大的矩阵代数内成为共轭。

\begin{corollary}\label{b2:cor:csa-field-normalizer}
设 \(k\) 为域，\(A\) 为中心单 \(k\) 代数，\(L\subseteq A\) 为含 \(k\) 的子域。令
\[
N_{A^\times}(L)=\{\,u\in A^\times\mid uLu^{-1}=L\,\}.
\]
共轭给出正合列
\[
1\longrightarrow C_A(L)^\times\longrightarrow N_{A^\times}(L)
\longrightarrow\operatorname{Aut}_k(L)\longrightarrow1.
\]
若 \([L:k]=\deg A\)，则核为 \(L^\times\)。
\end{corollary}
\begin{proof}
正规化条件保证共轭限制为 \(L\) 的 \(k\) 自同构。其核恰为在 \(L\) 上逐点恒等的那些可逆元，也就是 \(C_A(L)^\times\)。对任意 \(\sigma\in\operatorname{Aut}_k(L)\)，把包含 \(\iota:L\hookrightarrow A\) 与 \(\iota\circ\sigma\) 作为两种简单代数嵌入，\cref{b2:thm:skolem-noether}给出实现它的可逆元，因此该同态满射。最后的核识别由\cref{b2:prop:csa-self-centralizing-field}。
\end{proof}

满射保证对每个 \(\sigma\in\operatorname{Aut}_k(L)\)，存在 \(u_\sigma\in N_{A^\times}(L)\) 实现它，却不保证能同时选择这些元素，使 \(u_{\sigma\tau}=u_\sigma u_\tau\) 对全部 \(\sigma,\tau\) 成立。两边虽然实现同一个域自同构，比值仍可能是中心化子中的非平凡单位。\cref{b2:prop:quaternion-normalizer-nonsplit}将在实四元数中给出没有群截面的具体正合列；下一章的循环代数还将记录实现元相乘时的这种差异。

\subsection{适用条件与非单代数反例}
\label{b2:subsec:1703:03}

先看简单子代数假设。取 \(A=M_3(k)\)、\(B=k\times k\)，定义
\[
\varphi(x,y)=\operatorname{diag}(x,y,y),\qquad
\psi(x,y)=\operatorname{diag}(x,x,y).
\]
两者都是含幺的单射代数同态，但 \(\varphi(1,0)\) 的秩为一，\(\psi(1,0)\) 的秩为二。矩阵内共轭保持矩阵的秩，所以这两个嵌入不可能由内共轭相连。此例的环境 \(A\) 中心单，失败来自 \(B\) 有两个非零简单因子。对 \(\varphi\) 的表示，两种简单分量的重数为 \(1,2\)；对 \(\psi\) 则为 \(2,1\)。总维数同为三，分量重数却不同，正是前述证明只对单 Artin 环成立的原因。

\ProseParagraphBreak
中心性也不能任意删去。把 \(\mathbb C\) 看成实代数，它是单代数，但中心大于 \(\mathbb R\)。复共轭是非恒等实代数自同构，而交换代数中的每个内自同构都是恒等映射，所以复共轭并非内自同构。

\ProseParagraphBreak
非单代数甚至可以通过交换其不同因子产生外自同构。例如 \(k\times k\) 的映射 \((x,y)\mapsto(y,x)\) 保持代数结构，却不是内自同构，因为整个代数交换。因而“自同构都是内的”不是任意有限维代数的性质；它来自中心单性及前述简单模比较的共同作用。

% !TeX root = ../../Book2.tex
% 纲要标题：### 17.4 四元数代数、范数与分裂判据
% 本轮补充的纲要细目：
% * 以四元数左右乘构造范数相似变换，证明乘子及对角标量核；不由核的计算断言满射
% 纲要位置：计划纲要.md，第 1325 行。

\section{四元数代数、范数与分裂判据}
\label{b2:sec:1704}

四元数代数是次数二的中心单代数模型。它的乘法只需两个参数就能写下，是否分裂却与一个范数方程的可解性有关。本节假设 \(\operatorname{char}k\ne2\)，并固定 \(a,b\in k^\times\)；这一假设用于基的验证、共轭及二次扩张的描述。

% 纲要要点（供目录核对）：
%
% * 四元数代数
% * 矩阵代数及非分裂例子
% * 与二次型及范数方程的联系
%
% ---
%

\subsection{四元数代数}
\label{b2:subsec:1704:01}

\begin{definition}\label{b2:def:quaternion-algebra}
设 \(k\) 为特征不等于 \(2\) 的域，\(a,b\in k^\times\)。将自由结合代数 \(k\langle i,j\rangle\) 对下列关系生成的双边理想取商：
\[
i^2=a,\qquad j^2=b,\qquad ji=-ij
\]
所得含幺 \(k\) 代数称为\term[siyuanshu-daishu]{四元数代数}[quaternion algebra]，记为 \((a,b)_k\)。
\end{definition}
\symidx{\((a,b)_k\)：参数为 \(a,b\) 的四元数代数}

关系使每个字都能先把 \(i\) 移到 \(j\) 左边，再将两个指数降到零或一，所以 \(1,i,j,ij\) 生成这个代数。它们还线性无关，这一点不能仅由关系写法直接断言。

\begin{proposition}\label{b2:prop:quaternion-basis-splitting}
设 \(k\) 为特征不等于 \(2\) 的域，\(a,b\in k^\times\)，\(Q=(a,b)_k\)，则 \(Q\) 以 \(1,i,j,ij\) 为 \(k\) 基，并且是次数二的中心单 \(k\) 代数。若域扩张 \(K/k\) 含有 \(\alpha\) 满足 \(\alpha^2=a\)，则有分裂同构
\[
Q_K=K\otimes_kQ\longrightarrow M_2(K),\qquad
i\longmapsto I=\begin{pmatrix}\alpha&0\\0&-\alpha\end{pmatrix},
\quad j\longmapsto J=\begin{pmatrix}0&b\\1&0\end{pmatrix}.
\]
\end{proposition}
\begin{proof}
这两个矩阵满足 \(I^2=aI_2\)、\(J^2=bI_2\)、\(JI=-IJ\)，所以给出由商代数到矩阵代数的同态。在对角位置，\(I_2,I\) 线性无关，因为 \(2\alpha\ne0\)；在非对角位置，\(J,IJ\) 线性无关，因为其两个坐标组成的行列式为 \(-2\alpha b\ne0\)。故四个矩阵为 \(M_2(K)\) 的基。原四个生成元若有 \(k\) 线性关系，映射后也有同一关系，因此系数全零。这证明 \(\dim_kQ=4\)，扩张后同态便是基之间的同构。

现在验证 \(Q\) 本身中心单。若 \(U\) 为非零双边理想，则 \(K\otimes_kU\) 是 \(Q_K\cong M_2(K)\) 的非零理想，因矩阵环单而等于整个 \(Q_K\)。维数比较得 \(\dim_kU=4\)，所以 \(U=Q\)。若 \(z\in Z(Q)\)，其像在 \(M_2(K)\) 中为标量。以 \(I_2,I,J,IJ\) 为基比较系数，得到 \(z\) 只能是原 \(1\) 的 \(k\) 倍数，故中心为 \(k\)。维数为四使次数为二。
\end{proof}

对 \(q=x_0+x_1i+x_2j+x_3ij\)，定义它的\term[siyuanshu-gonge]{四元数共轭}[quaternion conjugation]为
\[
\bar q=x_0-x_1i-x_2j-x_3ij.
\]
生成元的像 \(i\mapsto-i,j\mapsto-j\) 在反代数中仍满足原关系，所以该映射反转乘法；它在上述基上作用两次为恒等，因此
\begin{equation}\label{b2:eq:quaternion-conjugation}
\overline{qr}=\bar r\bar q,\qquad\overline{\bar q}=q.
\end{equation}
直接按乘法关系展开，交叉项成对消去，得到
\begin{equation}\label{b2:eq:quaternion-norm}
q\bar q=\bar q q
=x_0^2-a x_1^2-b x_2^2+ab x_3^2.
\end{equation}

\begin{proposition}\label{b2:prop:quaternion-norm}
设 \(k\) 为特征不等于 \(2\) 的域，\(a,b\in k^\times\)，\(Q=(a,b)_k\)。对 \(q=x_0+x_1i+x_2j+x_3ij\in Q\)，令 \(\bar q=x_0-x_1i-x_2j-x_3ij\)。由于
\[
q\bar q=\bar q q=x_0^2-a x_1^2-b x_2^2+ab x_3^2\in k,
\]
可将这个标量记为 \(N(q)\)，称为\term[siyuanshu-fanshu]{四元数范数}[quaternion norm]。对任意 \(q,r\in Q\)，有
\[
N(qr)=N(q)N(r),\qquad
q\in Q^\times\Longleftrightarrow N(q)\ne0.
\]
可逆时 \(q^{-1}=N(q)^{-1}\bar q\)。在\cref{b2:prop:quaternion-basis-splitting}的分裂矩阵模型中，\(N(q)\) 就是矩阵行列式。
\end{proposition}
\begin{proof}
因 \(N(r)\in k\) 在中心，
\[
N(qr)=qr\bar r\bar q=q\cdot N(r)\bar q=N(q)N(r).
\]
若范数非零，\eqref{b2:eq:quaternion-norm}给出所述双侧逆。若 \(q\) 可逆，则乘法性给出 \(N(q)N(q^{-1})=N(1)=1\)，所以范数非零。

在分裂模型中，\(q\) 的矩阵为
\[
\begin{pmatrix}
x_0+\alpha x_1&b(x_2+\alpha x_3)\\
x_2-\alpha x_3&x_0-\alpha x_1
\end{pmatrix}.
\]
其行列式为 \(x_0^2-a x_1^2-bx_2^2+abx_3^2\)，正好是 \(N(q)\)。
\end{proof}
\symidx{\(N(q)=q\bar q\)：四元数范数}

这里的 \(N\) 是中心单代数约化范数在次数二情形的具体形式；\secref{b2:sec:1803}将从一般分裂矩阵模型重新得到这一识别。

\begin{proposition}\label{b2:prop:quaternion-norm-similitudes}
设 \(k\) 为特征不为二的域，\(Q=(a,b)_k\)、\(a,b\in k^\times\)，\(N(x)=x\bar x\) 为四元数范数。对 \(u,v\in Q^\times\)，定义 \(k\) 线性映射
\[
L_{u,v}:Q\longrightarrow Q,\qquad x\longmapsto uxv^{-1}.
\]
它是范数二次型的相似变换，乘子为 \(N(u)/N(v)\)。映射 \((u,v)\mapsto L_{u,v}\) 是从 \(Q^\times\times Q^\times\) 到该相似变换群的群同态，其核为 \(\{(c,c)\mid c\in k^\times\}\)。特别地，内自同构 \(x\mapsto uxu^{-1}\) 保持 \(N\)。
\end{proposition}
\begin{proof}
范数的对角系数为 \(1,-a,-b,ab\)，均非零，所以二次型非退化。\cref{b2:prop:quaternion-norm}的乘法性给出
\[
N(uxv^{-1})=\frac{N(u)}{N(v)}N(x).
\]
由特征不为二，极化同样使相应双线性型乘以这个标量。又有
\[
L_{u,v}L_{u',v'}(x)=uu'xv'^{-1}v^{-1}
=L_{uu',vv'}(x),
\]
故是群同态。若 \(L_{u,v}\) 为恒等，在 \(x=1\) 处即得 \(u=v\)；再对全部 \(x\) 使用等式，得到 \(u\) 属于 \(Q\) 的中心，也就是 \(k^\times\)。反过来，对角标量显然作用为恒等。这证明核的描述，而取 \(u=v\) 就得到等距结论。
\end{proof}

这给出一族由代数乘法产生的几何变换。仅从核的计算不能断言它们穷尽全部相似变换；这里只使用上述明确构造，也不预先调用正交群的进一步分类。

\subsection{矩阵代数与非分裂例子}
\label{b2:subsec:1704:02}

\begin{proposition}\label{b2:prop:quaternion-division-or-split}
设 \(k\) 为特征不等于 \(2\) 的域，\(a,b\in k^\times\)，\(Q=(a,b)_k\)，\(N:Q\rightarrow k\) 为其四元数范数，则 \(Q\) 或者是除代数，或者同构于 \(M_2(k)\)。它分裂，当且仅当二次型 \(N\) 有非平凡零点，即存在 \(q\in Q\setminus\{0\}\) 使 \(N(q)=0\)。
\end{proposition}
\begin{proof}
由\cref{b2:prop:csa-division-description}，\(Q\cong M_r(D)\)，其中 \(D\) 是除代数，且 \(4=r^2\dim_kD\)。若 \(r=1\)，\(Q\) 是除代数；否则只能 \(r=2,\dim_kD=1\)，即 \(D=k\)，所以 \(Q\cong M_2(k)\)。

由\cref{b2:prop:quaternion-norm}，每个非零元素可逆当且仅当每个非零元素的范数都非零。因此有非平凡零点恰好等价于不是除代数，按前一分类又等价于分裂。
\end{proof}

若 \(a=\alpha^2\in k^{\times2}\)，\cref{b2:prop:quaternion-basis-splitting}的两个矩阵已经定义在 \(k\) 中，因此 \((a,b)_k\) 直接分裂。同样，交换生成元 \(i,j\) 给出 \((a,b)_k\cong(b,a)_k\)，所以 \(b\) 为平方时也分裂。

\ProseParagraphBreak
实四元数代数
\[
\mathbb H=(-1,-1)_{\mathbb R}
\]
则不分裂，因为其范数为 \(x_0^2+x_1^2+x_2^2+x_3^2\)，对任意非零实向量严格为正。因而每个非零四元数都可逆；矩阵代数 \(M_2(\mathbb R)\) 却有非零奇异矩阵，二者不可能同构。扩张到 \(\mathbb C\) 后，参数 \(-1\) 成为平方，所以 \(\mathbb H\otimes_{\mathbb R}\mathbb C\cong M_2(\mathbb C)\)。

\ProseParagraphBreak
更一般地，对非零实数 \(a,b\)，只要其中一个为正，就因该参数为平方而分裂；若二者都为负，分别把 \(i,j\) 除以 \(\sqrt{-a},\sqrt{-b}\)，得到平方均为 \(-1\) 的生成元，故代数同构于 \(\mathbb H\)。因此实四元数代数恰有分裂矩阵代数与 \(\mathbb H\) 两种类型。这里分类的是四维四元数代数，没有在此断言所有实中心单代数的进一步分类。

\subsection{二次型与范数方程}
\label{b2:subsec:1704:03}

\begin{theorem}[四元数分裂的等价条件]\label{b2:thm:quaternion-splitting}
设 \(k\) 为特征不等于 \(2\) 的域，\(a,b\in k^\times\)，\(Q=(a,b)_k\)，其四元数范数为 \(N(q)=q\bar q\)。则
\[
Q\cong M_2(k)
\quad\Longleftrightarrow\quad Q\text{ 不是除代数}
\quad\Longleftrightarrow\quad
N\text{ 有非平凡零点}.
\]
若 \(a\) 不是 \(k\) 中的平方，令 \(L=k(\alpha)\)、\(\alpha^2=a\)，则上述条件还等价于
\[
b\in N_{L/k}(L^\times)
\quad\Longleftrightarrow\quad
b=x^2-a y^2\text{ 对某些 }x,y\in k.
\]
若 \(a\) 为非零平方，代数总分裂，最后一个方程也总有解。
\end{theorem}
\begin{proof}
\cref{b2:prop:quaternion-division-or-split}已经给出前三个条件的等价。以下处理 \(a\) 非平方时的域范数条件。

\(a\) 非平方时，\(L=k[i]\subseteq Q\) 是二次域，且 \(Q=L\oplus Lj\)。记 \(L\) 的非平凡自同构为 \(u\mapsto\bar u\)，则 \(ju=\bar u j\)。由前述范数公式，对 \(u,v\in L\) 有
\[
N(u+vj)=N_{L/k}(u)-b\cdot N_{L/k}(v).
\]
这里 \(N_{L/k}(x+y\alpha)=(x+y\alpha)(x-y\alpha)=x^2-a y^2\)，可由二次扩张的乘法矩阵或两个共轭直接计算。

若 \(Q\) 分裂，\cref{b2:prop:quaternion-division-or-split}给出非零 \(u+vj\) 的范数为零。不能有 \(v=0\)，否则域范数零使 \(u=0\)，矛盾。因此 \(v\ne0\)，得到
\(b=N_{L/k}(u)/N_{L/k}(v)=N_{L/k}(u/v)\)。反过来，若 \(b=N_{L/k}(z)\)，则 \(z+j\ne0\) 且其范数为零，所以代数分裂。域范数的显式公式又给出最后一个等价。

若 \(a=\alpha^2\) 且 \(\alpha\in k^\times\)，分裂性已证明，而
\[
x=\frac{b+1}{2},\qquad y=\frac{b-1}{2\alpha}
\]
满足 \(x^2-a y^2=b\)，故方程仍可解。
\end{proof}

已知范数方程的一个解，还可以直接构造分裂同构。若 \(a\) 非平方，取 \(z=x+y\alpha\in L\) 满足 \(N_{L/k}(z)=b\)。在二维 \(k\) 空间 \(L\) 上，让 \(i\) 作用为乘 \(\alpha\)，让 \(j\) 作用为 \(v\mapsto z\bar v\)。两者满足平方关系，且
\(z\overline{\alpha v}=-\alpha z\bar v\)，所以反交换，给出 \(Q\to\operatorname{End}_k(L)\)。它保持单位元，由单性为单射，两端维数同为四，故同构。在基 \(1,\alpha\) 下，具体矩阵为
\[
i\longmapsto\begin{pmatrix}0&a\\1&0\end{pmatrix},\qquad
j\longmapsto\begin{pmatrix}x&-ay\\y&-x\end{pmatrix}.
\]
后者的平方为 \((x^2-a y^2)I_2=bI_2\)，可作为公式的直接核验。

\ProseParagraphBreak
范数判据也等价于齐次方程
\[
X^2-aY^2-bZ^2=0
\]
有一个非零解。若 \(Z\ne0\)，除以 \(Z^2\) 就得到 \(b=(X/Z)^2-a(Y/Z)^2\)；若 \(Z=0\)，非零解必有 \(Y\ne0\)，于是 \(a=(X/Y)^2\)，同样推出分裂。反过来，范数方程解 \(b=x^2-a y^2\) 给出点 \((x,y,1)\)，平方参数时也可直接取 \((\alpha,1,0)\)。这样，四元数代数是否为矩阵代数，已经化为一个具体二次方程的可解性问题。

这里的范数既是二次型，也是除代数或矩阵代数的判别工具。若要比较它与 Clifford 代数、Witt 环中的不变量以及矩阵代数上的伴随对合，可接着读\secref{b2:sec:B02}和\secref{b2:sec:B03}；后者需要本章的中心单代数与 Skolem--Noether 定理。


\begin{chaptersummary}
中心单代数的正则双模是简单模，其双模自同态只有底域标量。稠密性因此把左右乘作用识别为全部线性自同态，得到包络代数同构。这个判据同时解释了任意标量扩张和张量积为何保持中心单性；在代数闭包上，结构定理又把代数化为全矩阵环，从而确定其次数。

\ProseParagraphBreak
一般简单子代数的中心可能大于底域，中心化子的中心恰好随之扩大。双中心化子和维数公式仍然成立；在共同中心 \(L\) 上取张量积，所得的是 \(C_A(L)\)，未必是整个 \(A\)。一般中心单代数中的包含极大子域未必有次数 \(\deg A\)；自中心化子域以及除代数中的最大子域则达到这个次数，并可分裂原代数。有限可分分裂域的证明通过非零多项式选择可分元素；约化特征多项式及其系数下降将在下一章建立。

\ProseParagraphBreak
Skolem–Noether 定理把简单代数的两种嵌入改写成两个同维模，由模的同构产生实现内共轭的可逆元。四元数代数则把这些抽象结构落实为四个基向量、四元数共轭和一个乘法范数：范数二次型有非平凡零点恰好等价于分裂，非平方参数情形又等价于二次扩张中的范数方程。左右乘以可逆四元数又给出范数的相似变换，内共轭则给出等距变换。下一章将用稳定等价、循环代数与约化范数进一步组织这些现象。
\end{chaptersummary}

\BookExercises

\begin{exercise}\label{b2:ex:17-simple-central-distinction}
设 \(k\) 为域、\(n\ge1\)，\(L/k\) 为非平凡有限扩张。分别说明 \(k[t]/(t^2)\)、\(L\)、\(M_n(k)\) 是否单、是否以 \(k\) 为中心。解释这些例子为何不能把定义中的两个要求合并为一个。
\end{exercise}
\begin{solution}
双数代数交换，中心是整个二维代数而非 \(k\)，且含非零真双边理想 \((\bar t)\)，所以既不单也不中心。域 \(L\) 是单环，但其中心为 \(L\ne k\)，所以不是中心单 \(k\) 代数。\(M_n(k)\) 单且中心为 \(kI_n\)，因此中心单。还可用 \(T_2(k)\) 分离另一方向：与 \(E_{11},E_{12}\) 交换的上三角矩阵只有标量矩阵，所以它的中心是 \(k\)，但严格上三角矩阵组成非零真双边理想，故它中心而不单。
\end{solution}

\begin{exercise}\label{b2:ex:17-inseparable-tensor}
设 \(k\) 为特征不等于二的域，\(d\in k^\times\) 非平方，\(L=k(\alpha)\)、\(\alpha^2=d\)。求 \(L\otimes_kL\) 的两个非平凡中心幂等元、全部理想及 Jacobson 根，说明它的非单性与\cref{b2:prop:inseparable-base-change-counterexample}中的非单性有何区别。
\end{exercise}
\begin{solution}
在 \(L\otimes_kL\) 中用第一个因子赋予左 \(L\) 结构，令 \(w=\alpha^{-1}\otimes\alpha\)，则 \(w^2=1\)，且 \(1,w\) 为 \(L\) 基。因此
\[
e_+=\frac{1+w}{2},\qquad e_-=\frac{1-w}{2}
\]
是非零正交幂等元，和为一，分别生成一维 \(L\) 因子。具体同构为
\[
L\otimes_kL\longrightarrow L\times L,\qquad
u\otimes v\longmapsto(uv,u\sigma(v)),
\]
其中 \(\sigma(\alpha)=-\alpha\)。它把 \(w\) 送到 \((1,-1)\)，故把 \(e_+,e_-\) 送到两个标准幂等元，保乘法并在基上为同构。

全部理想为 \(0,Le_+,Le_-\) 及全环，两极大理想的交为零，所以 Jacobson 根为零。这次非单来自分成两个域因子，没有非零幂零元；纯不可分反例则出现非零幂零元。两种情形的原代数 \(L\) 都不以 \(k\) 为中心。
\end{solution}

\begin{exercise}\label{b2:ex:17-quaternion-tensor-square}
设 \(k\) 为特征不等于二的域，\(a,b\in k^\times\)，\(Q=(a,b)_k\)。在 \(Q\otimes_kQ\) 中令
\[
s=-a^{-1}(i\otimes i),\qquad t=-b^{-1}(j\otimes j).
\]
证明 \(s,t\) 为交换的对合元，并求四个元素 \(e_{\epsilon,\delta}=(1+\epsilon s)(1+\delta t)/4\)、\(\epsilon,\delta\in\{1,-1\}\)，在四维空间 \(Q\) 上的作用。它们是否为本原幂等元？
\end{exercise}
\begin{solution}
由平方关系，\(s^2=t^2=1\)。两因子中的反交换同时产生两个负号，故 \(st=ts=(ab)^{-1}(ij\otimes ij)\)。由\cref{b2:prop:quaternion-tensor-square}，张量积作用为 \((q\otimes r)x=qx\bar r\)。于是 \(s\) 作用为 \(x\mapsto ixi^{-1}\)，在基 \(1,i,j,ij\) 上的符号为 \((+,+,-,-)\)；\(t\) 的符号为 \((+,-,+,-)\)。

因此 \(e_{+,+},e_{+,-},e_{-,+},e_{-,-}\) 分别投影到 \(k1,ki,kj,kij\)。它们非零、两两正交且和为一，矩阵模型中各有秩一。其角代数都是相应一维空间的自同态代数 \(k\)，没有非平凡幂等元，故四者都是本原幂等元。原四元数是否分裂不影响这个张量平方中的分解。
\end{solution}

\begin{exercise}\label{b2:ex:17-dimension-sixteen}
设 \(k\) 为域，\(A\) 为十六维中心单 \(k\) 代数。对 \(A\cong M_r(D)\)，列出可能的矩阵大小 \(r\) 及 \(\dim_kD\)。若 \(k\) 代数闭，或 \(k\) 有限，分别还能得到什么结论？
\end{exercise}
\begin{solution}
\(\deg A=4\)，而 \(\dim_kD\) 为平方，记 \(\deg D=d\)。维数关系为 \(16=r^2d^2\)，故 \(rd=4\)。可能的 \((r,\dim_kD)\) 为 \((1,16),(2,4),(4,1)\)。这只是参数的可能形式，不保证每个底域都有各类除代数。代数闭时有限维除代数只能是 \(k\)，所以 \(A\cong M_4(k)\)。有限底域时 \(D\) 本身也是有限除环，由\cref{b2:thm:wedderburn-little}交换，再由中心为 \(k\) 得 \(D=k\)，结论相同。
\end{solution}

\begin{exercise}\label{b2:ex:17-finite-field-maximal-subfield}
设 \(q\) 为素数幂、\(n\ge1\)。在 \(A=M_n(\mathbb F_q)\) 中构造一个次数为 \(n\) 的有限可分子域，并证明它的中心化子恰为自身。说明这个构造不需要在有限域上使用“非零多项式总有非零取值”的错误断言。
\end{exercise}
\begin{solution}
第一卷定理17.1.1保证存在有限域 \(L=\mathbb F_{q^n}\)，并可通过该卷定理17.2.3的子域描述包含 \(\mathbb F_q\)。选取 \(L\) 的 \(\mathbb F_q\) 基，左乘作用给出嵌入 \(L\hookrightarrow\operatorname{End}_{\mathbb F_q}(L)\cong M_n(\mathbb F_q)\)。若 \(T\) 与全部左乘交换，则 \(T(x)=x\cdot T(1)\)，故它也是乘以 \(L\) 中某个元素，中心化子就是 \(L\)。有限域扩张可分，因为每个元素满足无重根多项式 \(X^{q^n}-X\)。这里使用现成的有限域扩张与正则作用；例如 \(X^q-X\) 虽是非零多项式，却在 \(\mathbb F_q\) 每一点取零，不能作为无限域取点论证的类比。
\end{solution}

\begin{exercise}\label{b2:ex:17-tensor-factor-centralizer}
设 \(k\) 为域，\(A,B\) 为中心单 \(k\) 代数，在 \(C=A\otimes_kB\) 中把 \(A,B\) 识别为 \(A\otimes1,1\otimes B\)。证明它们互为在 \(C\) 内的中心化子。
\end{exercise}
\begin{solution}
这两个嵌入单射，因为 \(1\) 可以分别扩充为两个代数的基，张量基保证没有新线性关系。两因子逐元素交换，所以 \(1\otimes B\subseteq C_C(A\otimes1)\)。由\cref{b2:thm:csa-tensor-product}，环境 \(C\) 中心单；由\cref{b2:thm:csa-centralizer}，右端维数为 \(\dim_kC/\dim_kA=\dim_kB\)。故包含是相等，反方向交换两因子的角色即可。这里的中心化子在张量代数 \(C\) 内计算，而不是在其底向量空间的全部自同态中计算。
\end{solution}

\begin{exercise}\label{b2:ex:17-field-centralizer}
设 \(k\) 为域，\(L/k\) 为有限扩张、\([L:k]=d\)，\(V=L^r\)、\(r\ge1\)。将 \(L\) 通过标量乘法嵌入 \(A=\operatorname{End}_k(V)\)。求其中心化子和双中心化子，核对维数公式，且不假设 \(L/k\) 可分。
\end{exercise}
\begin{solution}
与全部 \(L\) 标量交换的 \(k\) 线性映射恰是 \(L\) 线性映射，因此中心化子为 \(\operatorname{End}_L(L^r)\cong M_r(L)\)。\(A\) 中心单，\(L\) 为简单子代数，\cref{b2:thm:csa-centralizer}给出双中心化子为原来的 \(L\)。维数分别是 \(\dim_kA=r^2d^2\)、\(\dim_kL=d\)、\(\dim_kM_r(L)=r^2d\)，确有 \(r^2d^2=d\cdot r^2d\)。
\end{solution}

\begin{exercise}\label{b2:ex:17-large-commutative-subalgebra}
设 \(k\) 为域，在 \(M_5(k)\) 中按 \(2+3\) 分块，令
\[
B=\left\{cI_5+\begin{pmatrix}0&X\\0&0\end{pmatrix}:c\in k, X\in M_{2\times3}(k)\right\}.
\]
计算 \(C_{M_5(k)}(B)\)、\(J(B)\) 及 \(B^\times\)，判断 \(B\) 是否为最大交换子代数，并比较其维数与环境代数的次数。
\end{exercise}
\begin{solution}
右上块的两个乘积为零，故 \(B\) 交换，维数为七。将一般矩阵写成 \(T=\begin{psmallmatrix}P&Q\\U&V\end{psmallmatrix}\)。与全部右上块交换等价于
\[
XU=0,\qquad UX=0,\qquad PX=XV
\quad\text{对所有 }X\in M_{2\times3}(k).
\]
用单条目矩阵 \(X=E_{ij}\) 检查，前两式迫使 \(U=0\)；第三式先迫使 \(P,V\) 都对角，再迫使所有对角条目为同一个 \(c\)。因此 \(T\in B\)，而 \(Q\) 任意，故中心化子就是 \(B\)。它因而最大交换。

令 \(I\) 为严格右上块组成的理想，则 \(I^2=0\)，且 \(B/I\cong k\)。幂零理想包含于根，商为域又使根不可能更大，故 \(J(B)=I\)。元素 \(cI_5+N_X\) 可逆当且仅当 \(c\ne0\)，此时逆为 \(c^{-1}I_5-c^{-2}N_X\)。其维数七大于 \(\deg M_5(k)=5\)，含非零平方零元，所以不是域。这是\cref{b2:prop:maximal-commutative-not-field}的另一个具体模型。
\end{solution}

\begin{exercise}\label{b2:ex:17-subfield-degree-divides}
令 \(D=(-1,-1)_{\mathbb Q}\)。比较二次域 \(K_+=\mathbb Q(\sqrt2)\) 与 \(K_-=\mathbb Q(\sqrt{-1})\)：它们能否嵌入 \(D\)，能否分裂 \(D\)？由此说明 \([K:\mathbb Q]\mid\deg D\) 能否作为嵌入的充分条件。
\end{exercise}
\begin{solution}
\(D\) 的范数为四个有理坐标的平方和，所以非零元素范数非零，\(D\) 是次数二的除代数。\(K_-\) 通过 \(\sqrt{-1}\mapsto i\) 嵌入 \(D\)，由\cref{b2:prop:quaternion-basis-splitting}也分裂 \(D\)。

若 \(K_+\) 嵌入 \(D\)，则某个 \(q=x_0+v\)、\(v=x_1i+x_2j+x_3ij\)，满足 \(q^2=2\)。其非标量部分为 \(2x_0v\)。若 \(v=0\)，就有有理数 \(x_0^2=2\)，不可能；若 \(v\ne0\)，则 \(x_0=0\)，而
\[
q^2=-(x_1^2+x_2^2+x_3^2)\le0,
\]
也不可能等于二。因此 \(K_+\) 不嵌入。将 \(K_+\) 视为实数子域后，同一个平方和范数在 \(K_+\) 上仍无非平凡零点，所以 \(D_{K_+}\) 仍为除代数，不分裂。两个域的次数都是二，都满足整除约束；\cref{b2:prop:csa-self-centralizing-field}给出的整除关系只是嵌入的必要条件。
\end{solution}

\begin{exercise}\label{b2:ex:17-inseparable-skolem}
设 \(k=\mathbb F_2(t)\)、\(L=k(\alpha)\)、\(\alpha^2=t\)。在 \(M_2(k)\) 中分别将 \(\alpha\) 送到
\[
I=\begin{pmatrix}0&t\\1&0\end{pmatrix},\qquad
J=\begin{pmatrix}0&1\\t&0\end{pmatrix}.
\]
证明这给出两个 \(k\) 代数嵌入，并找出实现共轭的可逆矩阵。
\end{exercise}
\begin{solution}
两矩阵平方都为 \(tI_2\)。\(t\) 不是 \(k\) 中的平方，故 \(X^2-t\) 不可约，两种代入都给出域 \(L\) 的含幺同态，因而单射。取 \(P=\operatorname{diag}(1,t)\)，直接计算得 \(PIP^{-1}=J\)。此扩张纯不可分，而两个嵌入仍共轭；这核验了 Skolem–Noether 定理不需要源代数中心可分的一个具体情形。
\end{solution}

\begin{exercise}\label{b2:ex:17-normalizer-nonsplit}
设 \(k\) 为特征不等于二的域，\(d\in k^\times\) 非平方，\(L=k(\sqrt d)\)。通过乘法作用把 \(L\) 嵌入 \(A=M_2(k)\)，其中
\[
\sqrt d\longmapsto T=\begin{pmatrix}0&d\\1&0\end{pmatrix}.
\]
求 \(N_{A^\times}(L)\)，并证明其到 \(\operatorname{Gal}(L/k)\) 的正合列具有群截面。与实四元数中的情形比较。
\end{exercise}
\begin{solution}
正则作用的中心化子为 \(L\)，或直接解与 \(T\) 交换的方程，得到矩阵 \(\begin{psmallmatrix}x&dy\\y&x\end{psmallmatrix}\)。其可逆元就是 \(L^\times\)。令 \(C=\operatorname{diag}(1,-1)\)，则 \(CTC^{-1}=-T\)、\(C^2=1\)，所以 \(C\) 实现二次域的非平凡自同构 \(\sigma\)。

任意正规化元限制到 \(L\) 后为恒等或 \(\sigma\)；前者属于 \(L^\times\)，后者与 \(C\) 的比值属于 \(L^\times\)。故
\[
N_{A^\times}(L)=L^\times\cup L^\times C
\cong L^\times\rtimes\operatorname{Gal}(L/k),
\]
其中 \(\sigma(z)=\bar z\)。映射 \(1\mapsto I_2\)、\(\sigma\mapsto C\) 是群截面。\cref{b2:prop:quaternion-normalizer-nonsplit}中所有实现复共轭的 \(zj\) 平方都为负实数；这里则确有平方为一的实现元，故两种正规化子扩张表现不同。
\end{solution}

\begin{exercise}\label{b2:ex:17-quaternion-parameter-change}
在 \(Q=(2,3)_{\mathbb Q}\) 中改取 \(i\) 与 \(j'=(1+i)j\) 为生成元。求新参数，并把得到的同构在两组四元数基中写成坐标变换。直接核对这个变换保留四元数范数。再判断 \((2,-6)_{\mathbb Q}\) 是否同构于 \(Q\)。
\end{exercise}
\begin{solution}
\((1+i)(1-i)=-1\)，由扭曲关系得
\[
(j')^2=(1+i)(1-i)j^2=-3,\qquad j'i=-ij'.
\]
又 \((1+i)^{-1}=i-1\)，所以 \(j=(i-1)j'\) 可恢复，给出 \((2,-3)_{\mathbb Q}\cong Q\)。把新基向量写回原基，有 \(j'=j+ij\)、\(ij'=2j+ij\)，故从新基坐标到原基坐标的变换为
\[
(x_0,x_1,x_2,x_3)\longmapsto
(x_0,x_1,x_2+2x_3,x_2+x_3).
\]
其矩阵行列式为负一，确为可逆变换。代入原范数，得到
\[
x_0^2-2x_1^2-3(x_2+2x_3)^2+6(x_2+x_3)^2
=x_0^2-2x_1^2+3x_2^2-6x_3^2,
\]
正是新参数 \((2,-3)\) 的范数。再把生成元 \(j\) 换成 \(ij\)，其平方为 \(-2\cdot3=-6\)，所以 \((2,-6)_{\mathbb Q}\) 也同构于 \(Q\)，符合\cref{b2:prop:quaternion-parameter-changes}。
\end{solution}

\begin{exercise}\label{b2:ex:17-rational-quaternions}
判断 \((2,3)_{\mathbb Q}\) 与 \((2,7)_{\mathbb Q}\) 哪一个分裂。对分裂者给出两生成元的 \(2\times2\) 有理矩阵；对不分裂者证明对应范数方程无有理解。
\end{exercise}
\begin{solution}
因为 \(7=3^2-2\cdot1^2\)，由\cref{b2:thm:quaternion-splitting}，\((2,7)_{\mathbb Q}\) 分裂，可取
\[
i\longmapsto\begin{pmatrix}0&2\\1&0\end{pmatrix},\qquad
j\longmapsto\begin{pmatrix}3&-2\\1&-3\end{pmatrix}.
\]
两矩阵平方分别为 \(2I_2,7I_2\)，并反交换，所以\cref{b2:prop:algebra-relations-universal}给出含幺同态。由\cref{b2:prop:quaternion-basis-splitting}，来源单且为四维，故这个非零同态单射，两边同维又使它成为同构。

若 \(x^2-2y^2=3\) 有有理解，清分母得到互素整数 \(u,v,w\)、\(w\ne0\)，满足 \(u^2-2v^2=3w^2\)。模三后为 \(u^2+v^2=0\)，只能 \(3\mid u,v\)，因为模三平方只有零和一。原等式继而使 \(3\mid w\)，与互素矛盾。因此该范数方程无解，由\cref{b2:thm:quaternion-splitting}，\((2,3)_{\mathbb Q}\) 不分裂；再由\cref{b2:prop:quaternion-division-or-split}，它为除代数。
\end{solution}

\begin{exercise}\label{b2:ex:17-quaternion-quadratic-centralizer}
设 \(k\) 为特征不等于二的域、\(a,b\in k^\times\)，在 \(Q=(a,b)_k\) 中求 \(C_Q\bigl(k[i]\bigr)\)，并证明 \(k[i]\) 是最大交换子代数。分别说明 \(a\) 非平方及 \(a=\alpha^2\) 时这个子代数的结构。
\end{exercise}
\begin{solution}
写 \(q=x_0+x_1i+x_2j+x_3ij\)，则
\(qi-iq=-2x_2ij-2a x_3j\)。基线性无关且 \(2a\ne0\)，所以交换当且仅当 \(x_2=x_3=0\)，即中心化子为 \(k+ki=k[i]\)。任何更大的交换子代数都落在这个中心化子内，故它最大交换。\(a\) 非平方时它是二次域；若 \(a=\alpha^2\)，则两个幂等元 \((1+i/\alpha)/2\)、\((1-i/\alpha)/2\) 非零、正交且和为一，给出 \(k[i]\cong k\times k\)。后者仍最大交换，却不是子域。
\end{solution}

\begin{exercise}\label{b2:ex:17-finite-quaternion-norm-count}
设 \(k=\mathbb F_q\)、\(q\) 为奇数，\(a,b\in k^\times\)。当 \(a\) 非平方时，证明方程 \(x^2-a y^2=b\) 恰有 \(q+1\) 组解；当 \(a\) 为平方时，恰有 \(q-1\) 组解。由此再证所有 \(k\) 上的四元数代数都分裂。
\end{exercise}
\begin{solution}
非平方时 \(L=\mathbb F_{q^2}=k(\sqrt a)\)，该方程是 \(N_{L/k}(x+y\sqrt a)=b\)。Frobenius 映射\footnote{弗罗贝尼乌斯（Ferdinand Georg Frobenius，1849-1917），德国数学家，研究群表示、矩阵和有限域。} \(z\mapsto z^q\) 固定 \(k\)，却不在全部 \(L\) 上为恒等，因为 \(X^q-X\) 至多有 \(q\) 个根。因此它是这个二次扩张的非平凡自同构，范数为 \(z\mapsto z\cdot z^q=z^{q+1}\)。第一卷推论17.2.1说明 \(L^\times\) 为阶 \(q^2-1\) 的循环群，因此这个映射核大小为 \(q+1\)，像大小为 \(q-1\)，正好满射到 \(k^\times\)。每个非零范数纤维都有 \(q+1\) 个元素，对应唯一的坐标对 \((x,y)\)。

若 \(a=\alpha^2\)，可逆线性换元 \(u=x+\alpha y,v=x-\alpha y\) 把方程化为 \(uv=b\)。每个 \(u\in k^\times\) 唯一确定 \(v=b/u\)，所以有 \(q-1\) 组解。两种情形都可解，\cref{b2:thm:quaternion-splitting}于是给出结论。
\end{solution}

\begin{exercise}\label{b2:ex:17-quaternion-automorphisms-norm}
设 \(k\) 为特征不等于二的域、\(a,b\in k^\times\)、\(Q=(a,b)_k\)。证明 \(Q\) 的每个 \(k\) 代数自同构都保留范数及共轭。说明这里不是单纯从某一组生成元的名字看出共轭不变性。
\end{exercise}
\begin{solution}
由\cref{b2:cor:csa-automorphisms-inner}，任一自同构为 \(q\mapsto uqu^{-1}\)。范数乘法性给出
\[
N(uqu^{-1})=N(u)N(q)N(u)^{-1}=N(q),
\]
所以范数被保持。记 \(T(q)=q+\bar q=2x_0\)。范数公式给出 \(T(q)=N(q+1)-N(q)-1\)，这是由代数的单位元及已经证明不变的范数决定的。因此 \(T\) 也不变，进而由 \(\bar q=T(q)-q\) 得共轭与自同构交换。证明没有要求自同构固定原来的 \(i,j\)。
\end{solution}

\begin{exercise}\label{b2:ex:17-quaternion-conjugation-rotation}
在实四元数 \(\mathbb H\) 中记 \(k_0=ij\)，并取 \(u=i+j\)。计算 \(x\mapsto uxu^{-1}\) 在基 \(1,i,j,k_0\) 上的作用，求其在三维实子空间 \(\mathbb Ri\oplus\mathbb Rj\oplus\mathbb Rk_0\) 上的行列式。
\end{exercise}
\begin{solution}
由 \(u^2=-2\)，有 \(u^{-1}=-(i+j)/2\)。按 \(i^2=j^2=-1\)、\(ji=-ij\) 相乘，得到 \(uiu^{-1}=j\)、\(uju^{-1}=i\)；该内自同构保持乘法，所以 \(uk_0u^{-1}=ji=-k_0\)，并且固定 \(1\)。三维子空间上的矩阵为
\[
\begin{pmatrix}0&1&0\\1&0&0\\0&0&-1\end{pmatrix},
\]
其行列式为一。\cref{b2:prop:quaternion-norm-similitudes}保证它保持范数的限制 \(x_1^2+x_2^2+x_3^2\)，这里可直接看出它固定 \(i+j\) 的方向，并把两个正交方向 \(i-j,k_0\) 变号。
\end{solution}

\begin{exercise}\label{b2:ex:17-explicit-quadratic-conjugator}
设 \(k\) 为特征不等于二的域、\(a,b\in k^\times\)，\(a\) 非平方、\(L=k(i)\subseteq(a,b)_k\)，\(i^2=a\)，并设 \(u\in L^\times\) 满足 \(N_{L/k}(u)=1\)。显式构造 \(z\in L^\times\)，使共轭 \(q\mapsto zqz^{-1}\) 固定 \(L\)，并将 \(j\) 送到 \(uj\)。
\end{exercise}
\begin{solution}
对 \(z\in L^\times\)，有 \(zjz^{-1}=z\bar z^{-1}j\)，所以只需求 \(z/\bar z=u\)。若 \(u\ne-1\)，取 \(z=1+u\ne0\)。由于 \(\bar u=u^{-1}\)，有 \(\bar z=1+u^{-1}=u^{-1}(1+u)\)，故 \(z/\bar z=u\)。若 \(u=-1\)，取 \(z=i\)，则 \(\bar z=-i\)，也有 \(z/\bar z=-1\)。\(L\) 交换，故上述共轭在 \(L\) 上恒等。这个构造给出 Skolem–Noether 共轭在二次域情形中的一个明确实现元。
\end{solution}

\begin{exercise}\label{b2:ex:17-nilpotent-maximal-commutative}
设 \(k\) 为域，在 \(A=M_4(k)\) 中令 \(N=\operatorname{diag}(J_2(0),J_2(0))\)。求 \(k[N]\) 的维数和中心化子。它是否最大交换？再求
\[
B=\operatorname{diag}\bigl(k[J_2(0)],k[J_2(0)]\bigr)
\]
的中心化子，解释单个 Jordan 块的条件为何重要。
\end{exercise}
\begin{solution}
\(N\ne0\)、\(N^2=0\)，故 \(k[N]\) 维数为二。把 \(k^4\) 视为 \((k[t]/(t^2))^2\)，\(N\) 就是乘 \(t\) 的算子。与它交换的算子恰为这个模的自同态，因而
\[
C_A(k[N])\cong M_2(k[t]/(t^2)),
\]
其 \(k\) 维数为八。例如 \(\operatorname{diag}(J_2(0),0)\) 与 \(N\) 交换，却不在 \(k[N]\) 中，添入它产生更大交换代数，所以 \(k[N]\) 不最大交换。

\(B\) 含两个块投影，故其中心化子必须分块对角；每个对角块又必须与 \(J_2(0)\) 交换，所以各块为 \(k[J_2(0)]\)。因此 \(C_A(B)=B\)，\(B\) 最大交换，维数为四，仍不是域。\cref{b2:prop:maximal-commutative-not-field}中的单个 Jordan 块有一个循环向量，此处两块的模有两份循环分量，允许额外的分量同态，故中心化子不再只是算子的多项式。
\end{solution}
